\documentclass[11pt,a4paper]{article}
% Préambule commun aux documents de PythonIPM/documentation.
% Compilation : tectonic <fichier>.tex
% tectonic compile avec XeTeX, qui lit l'UTF-8 nativement : surtout PAS inputenc ni fontenc,
% faits pour pdfLaTeX. Ce sont eux qui décodaient mal les caractères multi-octets — « cœur »
% ressortait en « cur », et les tirets cadratins en caractères parasites.
\usepackage{lmodern}
\usepackage[french]{babel}
\usepackage{microtype}
\usepackage[a4paper,margin=2.4cm]{geometry}
\usepackage{amsmath,amssymb}
\usepackage{booktabs}
\usepackage{longtable}
\usepackage{array}
\usepackage{ragged2e}
\usepackage{xcolor}
\usepackage{tcolorbox}
\usepackage{enumitem}
\usepackage{fancyhdr}
\usepackage[hidelinks]{hyperref}

\definecolor{encadre}{HTML}{F4F6F8}
\definecolor{bordure}{HTML}{4A6FA5}
\definecolor{alerte}{HTML}{B03A2E}

% Encadré « à retenir » : la lecture non technique du passage qui précède
\newtcolorbox{retenir}[1][]{
  colback=encadre, colframe=bordure, boxrule=0.6pt, arc=2pt,
  left=8pt, right=8pt, top=6pt, bottom=6pt,
  fonttitle=\bfseries, title={#1}
}

% Encadré d'avertissement : un point ouvert, un écart assumé
\newtcolorbox{attention}[1][]{
  colback=white, colframe=alerte, boxrule=0.6pt, arc=2pt,
  left=8pt, right=8pt, top=6pt, bottom=6pt,
  % le bandeau de titre a déjà la couleur d'alerte en fond : un titre `coltitle=alerte`
  % écrirait du rouge sur du rouge. On garde le blanc par défaut, comme l'encadré « retenir ».
  fonttitle=\bfseries, title={#1}
}

\newcommand{\col}[1]{\texttt{#1}}
% \ensuremath et non \textsubscript : \ipm est employé aussi bien dans le texte que dans les
% formules, et « M\textsubscript{0} » placé en mode mathématique est une erreur — silencieuse en
% PDF, mais qui fait échouer la conversion vers Word.
\newcommand{\ipm}{\ensuremath{M_{0}}}

\setlist[itemize]{itemsep=2pt, topsep=4pt}
\setlist[description]{style=nextline, leftmargin=1.2cm, itemsep=5pt}

% La source du document apparaît en tête de chaque page. Les documents RGPH la redéfinissent
% par \renewcommand{\source}{...} AVANT \begin{document}.
\newcommand{\source}{EHCVM 2021}

\pagestyle{fancy}
\fancyhf{}
\fancyhead[L]{\small\itshape IPM Côte d'Ivoire --- \source}
\fancyhead[R]{\small\thepage}
\renewcommand{\headrulewidth}{0.3pt}

\renewcommand{\arraystretch}{1.15}

\renewcommand{\source}{PND 2026-2030 / EHCVM vs RGPH}

\title{\bfseries EHCVM ou RGPH : quel IPM sert le mieux le PND 2026-2030 ?\\[4pt]
\large Analyse comparative et recommandation}
\author{Pipeline \col{PythonIPM}}
\date{Septembre 2026}

\begin{document}
\maketitle

\begin{abstract}
\noindent Deux indices de pauvreté multidimensionnelle ont été calculés pour la Côte d'Ivoire sur
l'année 2021 : l'un sur l'EHCVM, enquête auprès de 12 965 ménages, l'autre sur le RGPH, recensement
exhaustif de 5,6 millions de ménages. Ce document les confronte au Plan National de Développement
2026-2030 et tranche : lequel est le plus aligné, lequel aide le mieux, et comment les employer
ensemble. Il suppose connus les deux documents d'analyse qui le précèdent, mais peut se lire seul.
\end{abstract}

\begin{retenir}[La recommandation en une page]
\textbf{L'IPM EHCVM est le plus aligné sur le PND} : il couvre onze des vingt-deux cibles contre
sept, il est seul à mesurer les quatre cibles santé et nutrition du Pilier 4, seul à porter la
pauvreté monétaire, et seul à pouvoir être actualisé avant 2030. Si un seul indice devait être
retenu comme \emph{l'indice du PND}, c'est celui-là.

\textbf{Mais l'IPM RGPH est celui qui aide le mieux à agir} : il est seul à rendre opérationnels
trois instruments centraux du plan --- les Pôles Économiques Compétitifs, le registre social unifié
et le ciblage des 827 000 ménages de filets sociaux --- que l'enquête ne peut pas servir du tout.

\textbf{La recommandation n'est donc pas de choisir}, mais de répartir les rôles :
\begin{itemize}
  \item l'\textbf{EHCVM} tient le cadre de résultats --- la valeur 2021, la cible 2030, les points
        d'étape ;
  \item le \textbf{RGPH} tient la carte et le registre --- où investir, quels ménages servir ;
  \item la \textbf{variante harmonisée} fait le pont, et le fait bien : à indicateurs identiques,
        les deux sources donnent $\ipm = 0{,}2737$ et $0{,}2680$, soit 2~\% d'écart.
\end{itemize}
\end{retenir}

\tableofcontents
\bigskip

\section{Deux questions, pas une}

Le PND 2026-2030 pose au statisticien deux questions de nature différente, et l'essentiel de la
comparaison tient à ce qu'aucune source ne répond bien aux deux.

\begin{description}
  \item[Mesurer.] \emph{Où en est-on, et où en sera-t-on en 2030 ?} Le plan fixe des cibles
    chiffrées, pilier par pilier, avec une valeur de départ et une valeur d'arrivée. Y répondre
    demande un indicateur dont le contenu recoupe ces cibles, et qui puisse être recalculé en cours
    de plan.
  \item[Cibler.] \emph{Où faut-il investir, et quels ménages faut-il servir ?} Le Pilier 5 veut
    une « cartographie fine » pour ses Pôles Économiques Compétitifs, le Pilier 6 un « registre
    social unifié », le Pilier 4 une liste de 827 000 ménages. Y répondre demande une base
    exhaustive, ou au moins une maille territoriale très fine.
\end{description}

\begin{retenir}[Le partage naturel]
Une enquête par sondage est faite pour mesurer : questionnaire long, indicateurs riches, répétée
dans le temps. Un recensement est fait pour dénombrer : questionnaire court, indicateurs pauvres,
mais exhaustif et localisé. Le PND a besoin des deux, et la comparaison qui suit ne fait que le
chiffrer.
\end{retenir}

\section{Couverture des cibles du PND}

\subsection{Le décompte}

Vingt-deux cibles et orientations du PND peuvent être éclairées par un indice de pauvreté. Le
tableau ci-dessous en donne la couverture par source.

\begin{center}
\begin{tabular}{@{}lrr@{}}
\toprule
\textbf{Nature de la prise} & \textbf{EHCVM} & \textbf{RGPH} \\
\midrule
Directe & \textbf{11} & 7 \\
Approchée & 0 & 2 \\
Indirecte & 3 & 2 \\
Partielle & 2 & 0 \\
Opérationnelle (ciblage) & 0 & \textbf{3} \\
Aucune & 6 & 8 \\
\bottomrule
\end{tabular}
\end{center}

\subsection{Les cibles que l'une seule mesure}

C'est là que se joue l'essentiel. Onze cibles sont couvertes par une source et pas par
l'autre : sept par la seule EHCVM, quatre par le seul RGPH.

\begin{center}
\small
\begin{tabular}{@{}>{\RaggedRight}p{7.4cm}cc@{}}
\toprule
\textbf{Cible ou orientation du PND} & \textbf{EHCVM} & \textbf{RGPH} \\
\midrule
Enrôlés à la CMU de 20 à 30 millions (P.~4) & \textbf{oui} & non \\
Protection sociale de 27 à 50~\% (P.~4) & \textbf{oui} & non \\
Population à moins de 5 km d'un centre de santé, 82 à 100~\% (P.~4) & \textbf{oui} & non \\
Retard de croissance des moins de 5 ans, 21,4 à 18~\% (P.~4) & \textbf{oui} & non \\
Indice de l'écart genre, 0,655 à 0,7 (P.~4) & \textbf{oui} & non \\
Habitat décent : densité d'occupation (P.~5) & \textbf{oui} & non \\
Croisement avec la pauvreté monétaire (objectif général) & \textbf{oui} & non \\
\addlinespace
Espérance de vie de 62,3 à 65 ans (P.~4) & non & \textbf{oui} \\
Pôles Économiques Compétitifs, villes secondaires (P.~5) & non & \textbf{oui} \\
Registre social unifié pour le ciblage (P.~6) & non & \textbf{oui} \\
827 000 ménages bénéficiaires des filets sociaux (P.~4) & non & \textbf{oui} \\
\bottomrule
\end{tabular}
\end{center}

\begin{retenir}[Lecture]
L'EHCVM gagne sur le \textbf{contenu} du Pilier 4 : quatre cibles santé et nutrition, plus le genre,
que le recensement ne peut pas approcher, parce qu'il ne pose aucune question de santé.

Le RGPH gagne sur les \textbf{instruments} des Piliers 5 et 6 : la cartographie, le registre, le
ciblage --- trois orientations qui ne sont pas des cibles chiffrées mais des dispositifs à
construire, et qu'une enquête par sondage ne peut pas alimenter.
\end{retenir}

\section{La finesse territoriale : l'écart est massif}

Le Pilier 5 entend « une répartition optimale et équitable des investissements publics et privés »
fondée sur une « cartographie fine ». La question est donc : à quelle maille chaque source
produit-elle un chiffre publiable ?

Convention usuelle : une estimation dont le coefficient de variation dépasse 20~\% n'est pas
publiable.

\begin{center}
\begin{tabular}{@{}lrrrr@{}}
\toprule
& \multicolumn{2}{c}{\textbf{EHCVM}} & \multicolumn{2}{c}{\textbf{RGPH}} \\
\cmidrule(lr){2-3}\cmidrule(lr){4-5}
\textbf{Maille} & \textbf{Modalités} & \textbf{Non publ.} & \textbf{Modalités} &
  \textbf{Non publ.} \\
\midrule
Région & 33 & 0 & 33 & 0 \\
Département & 108 & 15 & 111 & \textbf{0} \\
Sous-préfecture ou commune & 442 & \textbf{92} & 519 & \textbf{19} \\
\bottomrule
\end{tabular}
\end{center}

\begin{attention}[Le chiffre qui tranche le Pilier 5]
Les 92 sous-préfectures non publiables de l'EHCVM représentent \textbf{47,9~\% de la population}
ivoirienne. Les 19 du RGPH en représentent \textbf{1,55~\%}.

Autrement dit : fonder une allocation d'investissements à la maille de la sous-préfecture sur
l'EHCVM reviendrait à ne rien pouvoir dire, de façon publiable, sur la moitié du pays. Le
recensement ramène cette zone d'ombre à un cinquantième.
\end{attention}

Le coefficient de variation médian donne la même lecture en continu : 15,4~\% contre 9,0~\% à la
sous-préfecture, 12,3~\% contre 5,0~\% au département.

\section{Le suivi dans le temps : l'écart est tout aussi massif, en sens inverse}

\begin{attention}[Le recensement ne dira pas si le PND a réussi]
Le RGPH est décennal : le prochain n'interviendra pas avant la fin du plan. L'IPM RGPH fournit une
valeur de départ, très détaillée, mais aucune évolution.

L'EHCVM est une enquête répétée --- 2018, 2021, et une vague dans l'intervalle du plan. C'est la
seule des deux sources capable d'alimenter en progrès le cadre de gouvernance que le PND met en
place.
\end{attention}

Pour un plan quinquennal assorti de cibles datées, c'est un argument qui ne se compense pas : un
indice qu'on ne peut recalculer qu'une fois n'est pas un indicateur de suivi.

\section{Ce que les deux mesurent également bien}

Sept indicateurs sont communs aux deux sources avec une définition quasi identique. Les taux de
privation se suivent de près, ce qui valide les deux calculs l'un par l'autre.

\begin{center}
\small
\begin{tabular}{@{}lrrl@{}}
\toprule
\textbf{Indicateur} & \textbf{EHCVM} & \textbf{RGPH} & \textbf{Commentaire} \\
\midrule
Année de scolarité & 62,9~\% & 63,2~\% & le diplôme approche bien les années d'études \\
Énergie de cuisson & 66,3~\% & 62,7~\% & --- \\
Emploi de subsistance & 40,5~\% & 46,0~\% & approximation côté RGPH \\
Fréquentation scolaire & 31,9~\% & 35,6~\% & --- \\
Biens d'équipement & 21,7~\% & 18,2~\% & le RGPH mesure les 8 biens, l'EHCVM 7 \\
Logement & 17,3~\% & 25,0~\% & nomenclatures de matériaux différentes \\
Électricité & 9,3~\% & 11,5~\% & --- \\
Chômage & 6,5~\% & 5,9~\% & --- \\
\addlinespace
Eau potable & 20,9~\% & 14,3~\% & le RGPH perd le temps de trajet \\
Toilette & 67,2~\% & 45,1~\% & le RGPH perd le partage des sanitaires \\
Alphabétisation & 62,9~\% & 70,4~\% & --- \\
État civil & 26,3~\% & 15,2~\% & --- \\
\bottomrule
\end{tabular}
\end{center}

\subsection{Le test de cohérence : la variante harmonisée}

Le pipeline produit, pour chaque source, une variante restreinte au socle des douze indicateurs
communs. C'est la seule comparaison légitime entre les deux.

\begin{center}
\begin{tabular}{@{}lrrr@{}}
\toprule
& $H$ & $A$ & \ipm \\
\midrule
EHCVM 2021 & 53,5~\% & 0,5113 & \textbf{0,2737} \\
RGPH 2021 & 52,0~\% & 0,5157 & \textbf{0,2680} \\
\bottomrule
\end{tabular}
\end{center}

\begin{retenir}[Un résultat qui compte pour le PND]
À indicateurs identiques, une enquête par sondage et un recensement exhaustif, collectés la même
année par deux dispositifs entièrement différents, donnent \textbf{2~\% d'écart}.

Conséquence pratique : la carte territoriale du recensement et le tableau de bord de l'enquête
peuvent être utilisés \emph{ensemble} sans craindre qu'ils racontent deux histoires différentes.
C'est ce qui rend l'architecture recommandée ci-dessous opérationnelle et non seulement
théorique.
\end{retenir}

\section{Le verdict}

\subsection{Lequel est le plus aligné ?}

\textbf{L'IPM EHCVM}, et pour trois raisons qui se cumulent.

\begin{enumerate}
  \item \textbf{Il couvre plus de cibles} : onze contre sept en prise directe, sur des cibles
        chiffrées assorties d'une valeur de départ et d'une valeur 2030.
  \item \textbf{Il couvre le cœur du Pilier 4} --- le pilier du capital humain, celui où se
        concentrent les cibles sociales du plan. Les quatre cibles santé et nutrition (CMU,
        protection sociale, accès aux soins, retard de croissance) sont hors de portée du
        recensement.
  \item \textbf{Il est actualisable}. Un cadre de résultats a besoin d'un indicateur qu'on peut
        recalculer ; le recensement ne le permettra pas avant 2031.
\end{enumerate}

\subsection{Lequel aide le mieux ?}

\textbf{Cela dépend de ce qu'on appelle aider}, et la réponse honnête est qu'ils aident à deux
moments différents du cycle de la politique publique.

\begin{center}
\small
\begin{tabular}{@{}>{\RaggedRight}p{4.6cm}>{\RaggedRight}p{4.6cm}>{\RaggedRight}p{4.6cm}@{}}
\toprule
\textbf{Moment} & \textbf{Ce que le PND demande} & \textbf{Source} \\
\midrule
Cadrer & fixer la cible de pauvreté à 2030, que le plan a laissée en blanc &
  EHCVM \\
Répartir & allouer les investissements du Pilier 5 entre territoires & \textbf{RGPH} \\
Cibler & constituer le registre social unifié, servir 827 000 ménages & \textbf{RGPH} \\
Suivre & mesurer les progrès à mi-parcours et en 2030 & EHCVM \\
Rendre compte & documenter l'atteinte des cibles du cadre de résultats & EHCVM \\
\bottomrule
\end{tabular}
\end{center}

\begin{attention}[Ce que choisir une seule source coûterait]
\textbf{Ne garder que l'EHCVM} : le Pilier 5 perdrait sa base statistique --- près de la moitié de
la population dans une zone dont l'IPM n'est pas publiable --- et le registre social unifié du
Pilier 6 devrait être construit sans information de privation.

\textbf{Ne garder que le RGPH} : le plan perdrait toute mesure de la santé, de la nutrition, de la
protection sociale et du genre, c'est-à-dire la plus grande partie des cibles sociales du Pilier 4,
et n'aurait aucun moyen de constater ses progrès avant 2031.
\end{attention}

\section{Architecture recommandée}

\subsection{Les trois rôles}

\begin{description}
  \item[L'EHCVM tient le cadre de résultats.] $\ipm = 0{,}2863$ en 2021 devient la valeur de
    départ ; une cible 2030 est fixée, comblant la seule absence de valeur chiffrée parmi les
    objectifs majeurs du plan ; onze cibles sectorielles sont suivies sur une série unique et
    cohérente plutôt que sur onze statistiques éparses.

  \item[Le RGPH tient la carte et le registre.] Un IPM publiable pour 111 départements et 500
    sous-préfectures fonde la planification spatiale du Pilier 5. Les 5,6 millions de ménages, dont
    1 751 746 pauvres et 285 115 en pauvreté sévère, chacun avec son score et sa localisation,
    fondent le registre social du Pilier 6 et le ciblage des filets sociaux.

  \item[La variante harmonisée fait le pont.] Elle garantit que les deux instruments parlent de la
    même chose. Elle sert aussi de contrôle : si un futur IPM EHCVM s'écartait fortement de la
    trajectoire, la comparaison harmonisée dirait s'il s'agit d'un effet réel ou d'un artefact de
    mesure.
\end{description}

\subsection{Ce qu'il reste à faire}

\begin{enumerate}
  \item \textbf{Faire adopter une cible d'IPM à 2030} dans le cadre de gouvernance du PND. C'est la
        seule action qui transforme un travail statistique en instrument de pilotage.
  \item \textbf{Publier la carte RGPH} à la maille du département et de la sous-préfecture, en
        signalant les 19 sous-préfectures non publiables.
  \item \textbf{Arbitrer les deux points méthodologiques ouverts} avant toute publication : le taux
        de privation de la toilette au RGPH, non réconcilié avec la production Stata antérieure, et
        l'approximation de l'emploi agricole de subsistance, qui porte à elle seule 26,8~\% de
        l'indice RGPH.
  \item \textbf{Reprendre le croisement monétaire / multidimensionnel} sur la méthode officielle de
        calcul du taux de pauvreté, pour documenter l'objectif général du plan.
  \item \textbf{Prévoir le recalcul de l'IPM EHCVM} sur la prochaine vague d'enquête, aux mêmes
        définitions --- le pipeline le permet sans modification.
\end{enumerate}

\section{Limites de cette comparaison}

\begin{enumerate}
  \item \textbf{Le décompte des cibles est un jugement, pas une mesure.} Classer une prise en
        « directe », « approchée » ou « indirecte » suppose d'apprécier la proximité entre un
        indicateur de ménage et une cible souvent formulée en taux individuel. Un autre analyste
        aboutirait à un décompte voisin, pas identique.
  \item \textbf{Le PND ne contient aucune cible d'IPM.} L'alignement établi ici est un alignement de
        contenu. Rien dans le plan n'engage à retenir cet indice.
  \item \textbf{Les deux IPM complets ne sont pas comparables en niveau}, la dimension Santé du
        RGPH étant presque inerte. Toutes les comparaisons de niveau de ce document passent par la
        variante harmonisée.
  \item \textbf{Les cibles du PND ont des années de référence hétérogènes} (2021, 2023, 2024, 2025)
        alors que les deux IPM sont datés de 2021. Les écarts à la cible ne sont donc pas des
        écarts à la valeur de départ officielle.
\end{enumerate}

\end{document}
